\documentclass[11pt]{article}

\usepackage[margin=1in]{geometry}
\usepackage[T1]{fontenc}
\usepackage[utf8]{inputenc}
\usepackage{lmodern}
\usepackage{amsmath}
\usepackage{amssymb}
\usepackage{booktabs}
\usepackage{hyperref}
\usepackage{xcolor}
\usepackage{listings}
\usepackage{enumitem}
\usepackage[numbers,sort&compress]{natbib}

\hypersetup{
    colorlinks=true,
    linkcolor=blue,
    urlcolor=blue,
    citecolor=blue
}

\lstdefinestyle{cppstyle}{
    language=C++,
    basicstyle=\ttfamily\small,
    keywordstyle=\color{blue!70!black},
    commentstyle=\color{green!40!black},
    stringstyle=\color{orange!70!black},
    breaklines=true,
    frame=single,
    columns=fullflexible,
    keepspaces=true
}

\lstdefinestyle{pythonstyle}{
    language=Python,
    basicstyle=\ttfamily\small,
    keywordstyle=\color{blue!70!black},
    commentstyle=\color{green!40!black},
    stringstyle=\color{orange!70!black},
    breaklines=true,
    frame=single,
    columns=fullflexible,
    keepspaces=true
}

\lstdefinestyle{shellstyle}{
    language=bash,
    basicstyle=\ttfamily\small,
    keywordstyle=\color{blue!70!black},
    commentstyle=\color{green!40!black},
    breaklines=true,
    frame=single,
    columns=fullflexible,
    keepspaces=true
}

\title{Automatic Differentiation in C++ Codes Using Order-Truncated Imaginary Numbers}
\author{
    Samuel Roberts\\
    The University of Texas at San Antonio
    \and
    Bruno Turcksin\\
    Oak Ridge National Laboratory
}
\date{\today}

\begin{document}

\maketitle

\begin{abstract}
Automatic differentiation is a useful tool for computing sensitivities, but
introducing derivative calculations into existing C++ simulation codes can
require substantial source-code modification. This paper presents
\texttt{cpp\_oti\_lib}, a header-only C++ implementation of order-truncated
imaginary (OTI) numbers for hypercomplex automatic differentiation. An OTI
algebra \(\mathbb{OTI}_m^n\) contains objects that propagate derivative
information with respect to \(m\) independent variables through total order
\(n\), allowing one overloaded model evaluation to return the function
value and the corresponding sensitivities. The library provides overloaded
arithmetic and elementary functions so that many numerical kernels can be
differentiated by changing selected scalar types. The implementation includes
focused unit tests and Python examples, and provides an optional
Kokkos-enabled path for CPU and GPU backends. Together, these features provide
a compact workflow for obtaining sensitivities from C++ simulation codes.
\end{abstract}

\tableofcontents

\section{Motivation and Significance}

\subsection{Automatic Differentiation}

Automatic differentiation (AD), also called algorithmic differentiation, is a
family of techniques for computing exact derivatives of functions represented
by computer programs \cite{Griewank2008,Baydin2015,Bucker2005}. AD computes
derivatives by systematically applying the chain rule to the elementary
operations and intrinsic functions that make up a program.

For example, consider a program written as a sequence of elementary operations.
Starting from an input \(v_0 = x\), the program computes intermediate values
\[
    v_i = \phi_i(v_{j_1}, v_{j_2}, \ldots),
    \qquad i = 1,\ldots,s,
\]
where each \(\phi_i\) is an elementary arithmetic operation or intrinsic
function and each dependency index satisfies \(j_k < i\). The final output is
\[
    y = f(x) = v_s
      = (\phi_s \circ \cdots \circ \phi_2 \circ \phi_1)(x),
\]
with the understanding that some \(\phi_i\) may take more than one argument.
AD augments this sequence so that derivative information is propagated along
with each computed value. If an intermediate variable is computed as
\[
    v = g(u),
\]
then its derivative is propagated locally through
\[
    \frac{dv}{dx} = \frac{dg}{du}\frac{du}{dx}.
\]
For programs with many intermediate variables, this local application of the
chain rule gives machine-precision derivatives without requiring the user to
manually differentiate the full program.

Two common AD modes are forward mode and reverse mode. Forward mode propagates
directional derivative information from inputs to outputs. It is especially
natural when the number of independent variables is modest or when derivatives
with respect to selected directions are needed. Reverse mode propagates
adjoint information from outputs back to inputs. It is especially efficient
for scalar-output functions with many input variables, such as objective
functions in optimization and machine learning \cite{Griewank2008,Baydin2015}.
AD tools and software systems have been developed for a wide range of
languages and application domains, including Fortran and C/C++ source
transformation tools, operator-overloading packages, robust optimization, and
modern machine learning systems
\cite{Bischof1996,Griewank1996,Hascoet2013,Su1997,Phipps2012,Abadi2016,Bradbury2018,Ansel2024}.

\subsection{Hypercomplex Algebras for Differentiation}

Hypercomplex differentiation methods extend real arithmetic by replacing a
real scalar with a structured object containing one real component and one or
more additional algebraic directions. Classical examples include complex
numbers, dual numbers, multicomplex numbers, multidual numbers, quaternions,
and other higher-dimensional algebras \cite{Kantor1989,Fischer2017}. In the
differentiation setting, these extra directions are not merely abstract
algebraic objects: they provide a way to propagate derivative information
through the same program that evaluates the original real-valued model.

The common computational pattern is to perturb an input variable along one or
more non-real directions, evaluate the target function using the algebra's
arithmetic rules, and then read derivatives from selected output
coefficients. This perspective is closely related to the complex-step method
for first derivatives and to dual-number automatic differentiation
\cite{Martins2003,Fischer2017}. The advantage is that derivative information
is carried by the algebra itself, which avoids subtractive cancellation in
finite differences and allows derivatives to be obtained to near machine
precision for analytic operations.

\subsection{Growth of Traditional Hypercomplex Representations}

Many hypercomplex derivative algebras are powerful but can become expensive as
the number of variables or derivative order increases. Multicomplex and
multidual constructions often grow by repeated doubling. This can be practical
for a small number of directions, but it becomes costly for multivariable
higher-order sensitivity analysis. For example, retaining all combinations
generated by a doubling process may store many coefficients that are not needed
for a particular total derivative order \cite{Fike2011,Lantoine2012,Cano2020}.

This scaling issue motivates the order-truncated imaginary algebra. Rather
than building a full recursively doubled algebra, OTI stores only the
multi-index terms whose total order is less than or equal to a prescribed
truncation order. If the goal is to compute all derivatives of a scalar model
with respect to \(M\) variables up to total order \(N\), the necessary number
of coefficients is
\[
    {M+N \choose N},
\]
which is exactly the coefficient count used by this library \cite{Cano2020}.

\subsection{Order-Truncated Imaginary Numbers}

Order-truncated imaginary numbers can be viewed as a sparse, derivative-driven
hypercomplex algebra. Each imaginary direction corresponds to a variable of
interest. Products of directions represent mixed derivative directions. Terms
whose total order exceeds the selected truncation order are discarded. In this
sense, an OTI number is naturally equivalent to a truncated multivariate
Taylor polynomial \cite{Cano2020}.

For a model evaluation
\[
    f(x_1,\ldots,x_M),
\]
the OTI inputs are initialized as
\[
    x_i \mapsto x_i + e_i.
\]
After the program is evaluated using OTI arithmetic, the output coefficients
contain the derivatives of \(f\) with respect to the selected variables. A
first-order coefficient corresponds to a gradient component, a second-order
coefficient corresponds to a Hessian entry after applying the appropriate
factorial scaling, and higher-order coefficients provide higher sensitivities.

The implementation in this repository focuses on the static scalar form of
this algebra. It does not attempt to reproduce every capability of larger OTI
or HYPAD libraries. Instead, it implements the core scalar arithmetic,
coefficient layout, elementary functions, and tests needed for a compact
header-only C++ library.

\subsection{Reference Context}

The background in this section follows the hypercomplex differentiation
viewpoint described in the HYPAD reference material provided with this
project, especially the discussion of hypercomplex algebras in Section 2.3.2.
The same general lineage includes complex-step differentiation
\cite{Martins2003}, hyper-dual numbers for second derivatives
\cite{Fike2011}, multicomplex methods for high-order derivatives
\cite{Lantoine2012}, and the OTI algebra developed for efficient
multivariable high-order derivative computation \cite{Cano2020}.

\subsection{Header-Only C++ OTI for Existing Codes}

The library implements scalar static OTI numbers as a compact C++ value type.
For practical purposes, an OTI number is a multivariate Taylor polynomial
truncated at a fixed total order. This makes the type useful for forward-mode
automatic differentiation, sensitivity analysis, local Taylor approximation,
and derivative-aware numerical algorithms.

Existing OTI implementations, such as the C-based \texttt{otilib}
library,\footnote{\texttt{otilib} project documentation:
\url{https://mauriaristi.github.io/otilib/}.} provide an important reference
for the algebra and its numerical behavior. However, introducing a C API into
an existing C++ code can require replacing arithmetic expressions with
library-specific function calls or manually changing operation definitions
throughout the source. The implementation described here uses a header-only C++
type and operator overloads so that many scalar kernels can be differentiated
by including the library and changing selected \texttt{double} variables to
\texttt{oti::otinum<M,N>}.

The current implementation is intentionally narrow in scope:

\begin{itemize}[noitemsep]
    \item Scalar OTI numbers are implemented.
    \item The number of variables and truncation order are compile-time
    template parameters.
    \item Coefficients are stored as \texttt{double}.
    \item The library is header-only for C++ users.
    \item A Python wrapper exists for testing and visualization, but the C++
    library remains the primary implementation.
\end{itemize}

Arrays, matrices, sparse variants, FEM integration, and quadrature utilities
are not part of the current implementation. They can be layered on later once
the scalar type is stable.

\section{Software Description}

\subsection{Mathematical Representation}

\subsubsection{Truncated Taylor Algebra}

An \texttt{oti::otinum<M,N>} object stores a truncated Taylor polynomial in
\texttt{M} formal infinitesimal variables,
\[
    e_1, e_2, \ldots, e_M,
\]
with all terms above total order \texttt{N} discarded. A generic element can
be written as
\[
    a = \sum_{|\alpha| \leq N} c_\alpha e^\alpha,
\]
where
\[
    \alpha = (\alpha_1,\ldots,\alpha_M), \qquad
    |\alpha| = \alpha_1 + \cdots + \alpha_M,
\]
and
\[
    e^\alpha = e_1^{\alpha_1} e_2^{\alpha_2} \cdots e_M^{\alpha_M}.
\]

The number of stored coefficients is
\[
    K(M,N) = {M+N \choose N}.
\]
For example, \texttt{otinum<2,2>} stores
\[
    {2+2 \choose 2} = 6
\]
coefficients.

\subsubsection{Coefficient Normalization}

The stored coefficients are normalized Taylor coefficients. For a scalar
function
\[
    f(x_1,\ldots,x_M),
\]
the coefficient associated with multi-index \(\alpha\) is
\[
    c_\alpha =
    \frac{1}{\alpha!}
    \frac{\partial^{|\alpha|} f}
         {\partial x_1^{\alpha_1}\cdots \partial x_M^{\alpha_M}},
\]
where
\[
    \alpha! = \alpha_1! \alpha_2! \cdots \alpha_M!.
\]

The C++ API exposes both the normalized coefficient and the ordinary partial
derivative:
\begin{itemize}[noitemsep]
    \item \texttt{coeff(alpha)} returns \(c_\alpha\).
    \item \texttt{partial(alpha)} returns \(\alpha! c_\alpha\).
\end{itemize}

This distinction is important. Raw coefficient access through
\texttt{operator[]} returns normalized coefficients, while \texttt{partial()}
returns derivative values in the usual calculus convention.

\subsubsection{Coefficient Layout}

Coefficients are stored in a flat \texttt{std::array<double, ncoeffs>}.
The implementation uses a graded multi-index layout: coefficients are grouped
by total order, and within each total order the exponent of earlier variables
is traversed from high to low.

For \texttt{otinum<2,2>}, the layout is:
\[
\begin{array}{ccl}
0 &:& (0,0) \\
1 &:& (1,0) \\
2 &:& (0,1) \\
3 &:& (2,0) \\
4 &:& (1,1) \\
5 &:& (0,2).
\end{array}
\]

This layout makes order truncation a prefix operation. All terms up to order
one are stored before all terms of order two, and so on.

\subsection{Repository Layout}

The current repository layout is:

\begin{lstlisting}[style=shellstyle]
include/otinum/otinum.hpp             # umbrella header
include/otinum/core.hpp               # otinum value type and arithmetic
include/otinum/functions.hpp          # public math functions
include/otinum/taylor.hpp             # scalar Taylor composition helpers
include/otinum/detail/binom.hpp       # constexpr binomial/factorial helpers
include/otinum/detail/multi_index.hpp # multi-index ranking and tables
examples/basic.cpp                    # simple C++ example
tests/smoke.cpp                       # broad smoke test
tests/test_*.cpp                      # focused unit tests
tests/run_unit_tests.sh               # focused test runner
bindings/python/otinum_py.cpp         # optional Python bindings
python_examples/*.py                  # Python visualization examples
logs/                                 # generated unit-test logs
\end{lstlisting}

\subsection{Implementation Architecture}

\subsubsection{Header Organization}

The public C++ interface is intentionally small. Most users include
\texttt{otinum/otinum.hpp}, which is the umbrella header for the scalar OTI
type and public mathematical functions. Internally, the implementation is
split into a few focused headers:

\begin{itemize}[noitemsep]
    \item \texttt{core.hpp} defines \texttt{oti::otinum<M,N>}, coefficient
    storage, construction, accessors, arithmetic, inversion, and truncated
    operations.
    \item \texttt{functions.hpp} defines the public \texttt{oti} namespace
    overloads for elementary functions.
    \item \texttt{taylor.hpp} implements the generic scalar Taylor-composition
    machinery used by \texttt{functions.hpp}.
    \item \texttt{detail/binom.hpp} provides the small compile-time
    combinatorial helpers used for coefficient counts and factorial scaling.
    \item \texttt{detail/multi\_index.hpp} defines the multi-index layout,
    ranking logic, order offsets, and multiplication tables.
    \item \texttt{detail/kokkos\_compat.hpp} centralizes the differences
    between standard C++ builds and Kokkos-enabled builds.
\end{itemize}

This organization keeps the user-facing type in one place while isolating the
indexing and backend-compatibility details that make the implementation work
for both ordinary host code and Kokkos device code.

\subsubsection{Static Storage and Compile-Time Tables}

\texttt{otinum<M,N>} is a fixed-size value type. It does not allocate dynamic
memory. Its coefficient array has exactly \({M+N \choose N}\) entries, and the
size is known at compile time. In a standard C++ build this array is a
\texttt{std::array}; in a Kokkos build it is a \texttt{Kokkos::Array}. The
choice is made through a shared alias in \texttt{detail/kokkos\_compat.hpp},
so the rest of the implementation can use the same source code in both modes.

The multi-index helper \texttt{detail::tables<M,N>} provides the metadata
needed to interpret the flat coefficient array. It stores or computes:

\begin{itemize}[noitemsep]
    \item the multi-index associated with each flat coefficient index,
    \item the total order of each coefficient,
    \item the starting flat index for each total order,
    \item the \(\alpha!\) scaling factor for each multi-index,
    \item the valid product triples used by truncated polynomial
    multiplication.
\end{itemize}

The multiplication table contains triples \((i,j,k)\) meaning that coefficient
\(i\) from the left operand and coefficient \(j\) from the right operand
contribute to output coefficient \(k\). Terms whose combined total order would
exceed \(N\) are omitted from the table. This means the runtime multiplication
loop does not need to rebuild multi-index sums or test every impossible
coefficient pair.

\subsubsection{Arithmetic Flow}

Linear operations are coefficient-wise. Adding two OTI numbers adds every
stored coefficient. Adding a scalar changes only coefficient zero, because the
scalar affects only the real component. Multiplying by a scalar scales every
coefficient, because all derivatives of the scaled function are scaled by the
same constant.

OTI multiplication is truncated polynomial convolution. The implementation
iterates over the precomputed product table and accumulates
\[
    out_k \mathrel{+}= lhs_i rhs_j.
\]
Because the product table contains only valid terms, all terms above the
configured total order are discarded automatically.

Division is implemented through an OTI inverse. For a value \(x\), the inverse
solves
\[
    x y = 1
\]
coefficient by coefficient. The real coefficient gives
\[
    y_0 = \frac{1}{x_0}.
\]
For every non-real coefficient, the target coefficient of \(xy\) is zero. The
graded coefficient layout ensures that when coefficient \(y_k\) is solved, the
remaining terms needed for that output coefficient have already been computed
or are known. The implementation therefore accumulates the known terms and
divides by the real part \(x_0\). This requires \(x_0 \neq 0\).

\subsubsection{Elementary Function Flow}

Elementary functions are implemented by composing a scalar Taylor series with
the nilpotent part of an OTI number. For an OTI value
\[
    x = x_0 + h,
\]
where \(x_0\) is the real coefficient and \(h\) has zero real coefficient, a
scalar function \(g\) is evaluated as
\[
    g(x_0 + h)
    =
    \sum_{k=0}^{N}
    \frac{g^{(k)}(x_0)}{k!} h^k.
\]
The series is finite because \(h\) is nilpotent under total-order truncation.
The helper \texttt{detail::apply\_scalar} implements this pattern. Function
specific helpers such as \texttt{exp\_coeffs}, \texttt{log\_coeffs},
\texttt{pow\_coeffs}, \texttt{sin\_coeffs}, and \texttt{cos\_coeffs} provide
the scalar coefficients \(g^{(k)}(x_0)/k!\).

Derived functions are intentionally expressed through simpler primitives where
that keeps the implementation consistent. For example, \texttt{sqrt(x)} uses
\texttt{pow(x, 0.5)}, \texttt{tan(x)} uses \texttt{sin(x) / cos(x)}, and
hyperbolic functions are composed from \texttt{exp}. This reduces the number
of independent derivative formulas that must be maintained.

\subsubsection{Kokkos Compatibility Layer}

Kokkos support is enabled with \texttt{OTI\_ENABLE\_KOKKOS}. In this mode,
\texttt{detail/kokkos\_compat.hpp} switches fixed-size storage from
\texttt{std::array} to \texttt{Kokkos::Array}, routes scalar math calls through
Kokkos math functions, and defines the annotation macros used by the rest of
the library. The most important macros are:

\begin{center}
\begin{tabular}{lp{0.55\linewidth}}
\toprule
Macro & Meaning \\
\midrule
\texttt{OTI\_FUNCTION} & Function callable in the active build mode \\
\texttt{OTI\_CONSTEXPR\_FUNCTION} & \texttt{constexpr} plus device annotation when needed \\
\texttt{OTI\_ASSERT} & Host or Kokkos assertion depending on build mode \\
\bottomrule
\end{tabular}
\end{center}

This layer allows the same arithmetic source to be instantiated in ordinary
C++ code and inside Kokkos kernels.

\subsection{Using the C++ Library}

\subsubsection{Including the Header}

The library is header-only. A C++ user only needs the include directory:

\begin{lstlisting}[style=cppstyle]
#include "otinum/otinum.hpp"
\end{lstlisting}

No linking step is required.

\subsubsection{Compiling a Program}

The basic example can be compiled directly:

\begin{lstlisting}[style=shellstyle]
c++ -std=c++17 -I include examples/basic.cpp -o /tmp/otinum_basic
/tmp/otinum_basic
\end{lstlisting}

\subsubsection{Basic Example}

\begin{lstlisting}[style=cppstyle]
#include <iostream>

#include "otinum/otinum.hpp"

int main()
{
    using T = oti::otinum<2, 2>;

    T x = T::variable(0, 1.5);
    T y = T::variable(1, 0.3);
    T f = oti::sin(x * y) + oti::exp(x);

    std::cout << "f = " << f.real() << '\n';
    std::cout << "df/dx = " << f.partial({1, 0}) << '\n';
    std::cout << "df/dy = " << f.partial({0, 1}) << '\n';
    std::cout << "d2f/dxdy = " << f.partial({1, 1}) << '\n';
}
\end{lstlisting}

\subsubsection{Constructing Variables}

\texttt{T::variable(i, value)} creates the OTI number
\[
    value + e_i.
\]
For example:

\begin{lstlisting}[style=cppstyle]
using T = oti::otinum<3, 2>;

T x1 = T::variable(0, 1.0);
T x2 = T::variable(1, 2.0);
T x3 = T::variable(2, 3.0);
\end{lstlisting}

The variable index is zero-based in C++. For a three-variable type,
\texttt{partial({1,0,0})} corresponds to differentiation with respect to the
first variable, \texttt{partial({0,1,0})} to the second, and
\texttt{partial({0,0,1})} to the third.

\subsubsection{Direct Coefficient Construction}

The library also supports construction from a complete coefficient array:

\begin{lstlisting}[style=cppstyle]
using T = oti::otinum<2, 2>;

std::array<double, T::ncoeffs> coeffs{};
coeffs[0] = 1.0; // real part
coeffs[1] = 2.0; // normalized coefficient for (1,0)
coeffs[2] = 3.0; // normalized coefficient for (0,1)

T value = T::from_coeffs(coeffs);
\end{lstlisting}

\subsubsection{Accessing Values and Derivatives}

Important accessors include:

\begin{center}
\begin{tabular}{ll}
\toprule
Method & Meaning \\
\midrule
\texttt{real()} & Real coefficient, index zero \\
\texttt{operator[](i)} & Raw normalized coefficient by flat index \\
\texttt{coeff(alpha)} & Normalized coefficient for a multi-index \\
\texttt{partial(alpha)} & Ordinary partial derivative value \\
\texttt{data()} & Const reference to the coefficient array \\
\bottomrule
\end{tabular}
\end{center}

Multi-indices outside the configured total order return zero from
\texttt{coeff()} and \texttt{partial()}.

\subsection{Implemented Operations}

\subsubsection{Arithmetic}

The core arithmetic operators are implemented for \texttt{otinum} values and
for combinations with \texttt{double}:

\begin{itemize}[noitemsep]
    \item Addition: \texttt{+}
    \item Subtraction: \texttt{-}
    \item Unary negation: \texttt{-x}
    \item Multiplication: \texttt{*}
    \item Division: \texttt{/}
\end{itemize}

Addition, subtraction, unary negation, and scalar multiplication are
coefficient-wise operations. Multiplication is a truncated polynomial
convolution:
\[
    c_{\alpha+\beta} \mathrel{+}= a_\alpha b_\beta
    \qquad \text{when } |\alpha+\beta| \leq N.
\]

Division is implemented as multiplication by the inverse. If
\[
    x = r(1+h),
\]
where \(r\) is the real part and \(h\) has zero real part, then \(h\) is
nilpotent under total-order truncation. Therefore,
\[
    \frac{1}{x}
    =
    \frac{1}{r}
    \sum_{k=0}^{N} (-h)^k.
\]

\subsubsection{Truncated Operations}

The following helper operations are implemented:

\begin{itemize}[noitemsep]
    \item \texttt{trunc\_mul(a,b,max\_order)}
    \item \texttt{trunc\_add(a,b,max\_order)}
    \item \texttt{gem(a,b,c)}, equivalent to \(a b + c\)
\end{itemize}

These are useful when a caller explicitly wants to discard terms above a
specified order.

\subsubsection{Math Functions}

The public math functions currently include:

\begin{center}
\begin{tabular}{llll}
\toprule
\texttt{exp} & \texttt{log} & \texttt{log10} & \texttt{logb} \\
\texttt{pow} & \texttt{sqrt} & \texttt{cbrt} & \texttt{abs} \\
\texttt{sin} & \texttt{cos} & \texttt{tan} & \\
\texttt{sinh} & \texttt{cosh} & \texttt{tanh} & \\
\bottomrule
\end{tabular}
\end{center}

The implementation uses univariate Taylor composition. For a scalar function
\(g\) applied to an OTI number \(x\), write
\[
    x = x_0 + h,
\]
where \(x_0 = \texttt{x.real()}\) and \(h\) has zero real part. Since \(h\) is
nilpotent,
\[
    g(x) =
    \sum_{k=0}^{N}
    \frac{g^{(k)}(x_0)}{k!} h^k.
\]

The helper \texttt{detail::apply\_scalar} implements this finite expansion.
Function-specific helpers provide the scalar Taylor coefficients
\(g^{(k)}(x_0)/k!\).

\subsubsection{Domain and Differentiability Behavior}

The elementary function overloads follow the real-valued domain behavior of
the underlying scalar function evaluated at \texttt{value.real()}. For
example, \texttt{log(x)} requires a positive real expansion point for
real-valued derivatives, \texttt{inv(x)} and division require a nonzero real
part in the denominator, and \texttt{tan(x)} is singular where
\(\cos(x_0)=0\). Non-integer powers can require positive or nonzero expansion
points depending on the exponent and derivative order.

The \texttt{abs} overload is piecewise. Away from zero it returns either
\(x\) or \(-x\), matching the derivative of the active branch. At real part
zero, the mathematical derivative is not defined; the implementation follows
its branch rule rather than attempting to define a generalized derivative.

\subsection{Python Wrapper}

\subsubsection{Purpose}

The Python wrapper is optional and is intended for testing, visualization, and
interactive exploration. It does not change the C++ implementation. The C++
library remains header-only and can still be consumed without Python,
\texttt{pybind11}, CMake, or packaging tools.

\subsubsection{Build and Installation}

The Python extension is built with \texttt{pybind11},
\texttt{scikit-build-core}, and CMake. From the repository root:

\begin{lstlisting}[style=shellstyle]
python -m pip install -e .
\end{lstlisting}

The package can then be imported:

\begin{lstlisting}[style=pythonstyle]
import otinum as oti
\end{lstlisting}

\subsubsection{Exposed Types}

The wrapper exposes a fixed grid of small concrete template instantiations:

\begin{center}
\begin{tabular}{ccc}
\toprule
\texttt{OTI\_1\_1} & \texttt{OTI\_1\_2} & \texttt{OTI\_1\_3} \\
\texttt{OTI\_2\_1} & \texttt{OTI\_2\_2} & \texttt{OTI\_2\_3} \\
\texttt{OTI\_3\_1} & \texttt{OTI\_3\_2} & \texttt{OTI\_3\_3} \\
\bottomrule
\end{tabular}
\end{center}

These cover the current visualization and testing use cases while keeping the
compiled extension small. Larger or additional \((M,N)\) combinations can be
bound later if needed.

\subsubsection{Python Usage Example}

\begin{lstlisting}[style=pythonstyle]
import otinum as oti

T = oti.OTI_2_3

x = T.variable(0, 1.5)
y = T.variable(1, 0.3)

f = oti.sin(x * y) + oti.exp(x)

print(f.real())
print(f.partial([1, 0]))
print(f.partial([0, 1]))
print(f.data())
\end{lstlisting}

\subsubsection{Python API Behavior}

The Python wrapper supports:

\begin{itemize}[noitemsep]
    \item Default and real-valued constructors.
    \item \texttt{variable(i, value=0.0)} with bounds checking.
    \item \texttt{from\_coeffs(coeffs)}.
    \item \texttt{real()}, \texttt{data()}, \texttt{coeff(alpha)}, and
    \texttt{partial(alpha)}.
    \item Python coefficient indexing through \texttt{\_\_getitem\_\_} and
    \texttt{\_\_setitem\_\_}.
    \item Same-type arithmetic and scalar arithmetic.
    \item Module-level math functions such as \texttt{oti.sin(x)} and
    \texttt{oti.exp(x)}.
\end{itemize}

Different OTI classes do not mix implicitly. For example, adding
\texttt{OTI\_2\_2} to \texttt{OTI\_2\_3} raises a Python \texttt{TypeError}.
This avoids ambiguous promotion rules between different truncated algebras.

\subsection{Kokkos-Enabled Performance-Portable Backend}

The software also provides a Kokkos-enabled form of the scalar static OTI
type. This backend preserves the same algebraic interface while making the
core arithmetic callable from Kokkos execution spaces. The coefficient storage
is represented with \texttt{Kokkos::Array} rather than \texttt{std::array}, and
the constructors, arithmetic operators, coefficient accessors, and elementary
function kernels are annotated with \texttt{KOKKOS\_FUNCTION} where device
execution is required.

This design keeps the user-facing programming model close to the CPU-only
header implementation. A scalar kernel written in terms of overloaded
arithmetic can be instantiated with the OTI type and executed inside a Kokkos
parallel loop. The result is that OTI derivative propagation is lifted from a
single CPU execution path to performance-portable CPU and GPU backends without
changing the mathematical API.

The Kokkos tests mirror the focused CPU unit tests. They check coefficient
layout, arithmetic, elementary functions, truncation behavior, and derivative
extraction on Kokkos-supported execution spaces. This ensures that the GPU
execution path preserves the same derivative semantics as the standard
header-only CPU implementation.

\section{Illustrative Examples}

The repository includes several Python examples under \texttt{python\_examples}.
Each script writes generated PDF figures into:

\begin{lstlisting}[style=shellstyle]
python_examples/figures/<script_name>/
\end{lstlisting}

\subsection{One-Dimensional Example}

\texttt{python\_examples/one\_dimensional.py} uses \texttt{OTI\_1\_3} to
visualize a scalar one-dimensional function and its first three derivatives.
It also compares the true function to a third-order Taylor approximation and
plots the Taylor approximation error.

Generated figures include:

\begin{itemize}[noitemsep]
    \item \texttt{function\_and\_derivatives.pdf}
    \item \texttt{third\_order\_taylor\_comparison.pdf}
    \item \texttt{third\_order\_taylor\_error.pdf}
\end{itemize}

\subsection{Newton-Raphson Example}

\texttt{python\_examples/1d\_newton\_raphson.py} solves
\[
    F(u,a,b) = u^3 + a u - b = 0
\]
with Newton-Raphson, where \(a\) and \(b\) are OTI design parameters. The
script tracks residual convergence, Newton step size, and changes in
derivative coefficients. This gives a derivative-aware stopping criterion for
the numerical method.

Generated figures include:

\begin{itemize}[noitemsep]
    \item \texttt{newton\_convergence\_metrics.pdf}
    \item \texttt{derivative\_history.pdf}
    \item \texttt{design\_parameter\_taylor\_slice.pdf}
\end{itemize}

\subsection{Two-Dimensional Example}

\texttt{python\_examples/two\_dimensional.py} uses \texttt{OTI\_2\_3} to
evaluate a two-variable scalar function, gradient components, a Hessian
determinant, and a third-order Taylor approximation error over a grid.

Generated figures include:

\begin{itemize}[noitemsep]
    \item \texttt{surface.pdf}
    \item \texttt{contours.pdf}
    \item \texttt{gradient\_field.pdf}
    \item \texttt{hessian\_determinant.pdf}
    \item \texttt{third\_order\_taylor\_error.pdf}
\end{itemize}

\subsection{Three-Dimensional Example}

\texttt{python\_examples/three\_dimensional.py} uses \texttt{OTI\_3\_2} to
evaluate a scalar field, gradient norm, Laplacian, a high-value region
visualization, and a second-order Taylor approximation error on a slice.

Generated figures include:

\begin{itemize}[noitemsep]
    \item \texttt{field\_slice.pdf}
    \item \texttt{gradient\_norm\_slice.pdf}
    \item \texttt{laplacian\_slice.pdf}
    \item \texttt{high\_value\_region.pdf}
    \item \texttt{second\_order\_taylor\_error\_slice.pdf}
\end{itemize}

\section{Benchmarks and Verification}

\subsection{Motivation}

The original \texttt{tests/smoke.cpp} file performs many checks in one
executable. That is useful as a broad regression test, but it is less
convenient for debugging because a failure may come from many possible
subsystems.

The focused unit test suite splits that coverage into smaller executables.
Each file targets one library area, so failures are easier to localize.

\subsection{Test Files}

\begin{center}
\begin{tabular}{ll}
\toprule
Test file & Purpose \\
\midrule
\texttt{test\_layout\_tables.cpp} & Binomial counts, ranks, order offsets, table consistency \\
\texttt{test\_construction\_access.cpp} & Constructors, variables, raw access, \texttt{deriv}, \texttt{partial} \\
\texttt{test\_linear\_arithmetic.cpp} & Addition, subtraction, scalar operations, unary negation \\
\texttt{test\_multiplication\_division.cpp} & Polynomial products, powers by multiplication, inverse/division \\
\texttt{test\_truncated\_operations.cpp} & \texttt{trunc\_mul}, \texttt{trunc\_add}, and \texttt{gem} \\
\texttt{test\_exp\_log\_pow.cpp} & Exponential, logarithm, power, square root, log bases \\
\texttt{test\_trig\_hyperbolic.cpp} & Trigonometric and hyperbolic identities \\
\texttt{test\_abs\_large\_shapes.cpp} & \texttt{abs} behavior and selected larger template shapes \\
\bottomrule
\end{tabular}
\end{center}

\subsection{Running the Tests}

The focused unit tests are run with:

\begin{lstlisting}[style=shellstyle]
tests/run_unit_tests.sh
\end{lstlisting}

The runner compiles each \texttt{tests/test\_*.cpp} file into a separate
executable under:

\begin{lstlisting}[style=shellstyle]
/tmp/otinum_unit_tests
\end{lstlisting}

Then it runs each executable in sequence.

\subsection{Test Logs}

Each test run writes a timestamped log file under \texttt{logs/}:

\begin{lstlisting}[style=shellstyle]
logs/unit_tests_YYYYMMDD_HHMMSS.log
\end{lstlisting}

The log includes:

\begin{itemize}[noitemsep]
    \item timestamp,
    \item compiler,
    \item compiler flags,
    \item build directory,
    \item each compile command,
    \item each executable invocation,
    \item output from passing tests,
    \item final suite result.
\end{itemize}

This is useful for reporting because passing test output is preserved, not
only failures.

Output locations can be overridden:

\begin{lstlisting}[style=shellstyle]
BUILD_DIR=/tmp/my_otinum_tests LOG_DIR=/tmp/my_otinum_logs tests/run_unit_tests.sh
\end{lstlisting}

\subsection{Original Smoke Test}

The original smoke test remains available:

\begin{lstlisting}[style=shellstyle]
c++ -std=c++17 -I include tests/smoke.cpp -o /tmp/otinum_smoke
/tmp/otinum_smoke
\end{lstlisting}

It is still useful as a broad all-in-one regression check.

\subsection{Optional Operations Microbenchmark}

The repository includes an optional CPU microbenchmark that reports timings
for addition, multiplication, division, exponential, sine, and a mixed
expression across several \texttt{oti::otinum<M,N>} shapes. Build and run it
separately from the correctness tests:

\begin{lstlisting}[style=shellstyle]
cmake -S . -B build-benchmarks \
  -DOTI_BUILD_TESTS=OFF \
  -DOTI_BUILD_BENCHMARKS=ON
cmake --build build-benchmarks --parallel
./build-benchmarks/benchmark_operations
\end{lstlisting}

The executable writes CSV rows containing elapsed seconds and algebra-size
metadata. These measurements are intended for comparisons made on the same
machine and compiler configuration. The repository does not currently provide
a portable CPU/GPU benchmark suite or publish cross-platform performance
claims.

\section{Impact}

The main impact of the software is a lower adoption barrier for higher-order
automatic differentiation in existing C++ simulation codes. By using a
header-only implementation and C++ operator overloading, users can often add
derivative propagation by changing selected scalar types while preserving the
structure of the original numerical kernel. This is particularly useful for
engineering codes where the cost of manually deriving and maintaining
gradient, Hessian, or higher-order sensitivity expressions is prohibitive.

The static OTI representation also addresses a practical storage issue in
hypercomplex derivative methods. By storing only multi-index coefficients up
to the requested total order, the library avoids propagating unnecessary
directions. This makes the implementation suitable for scalar derivative
calculations where the desired number of variables and derivative order are
known in advance.

The Kokkos-enabled backend extends this impact to performance-portable
computing. Many scientific applications evaluate the same scalar kernel over
large collections of elements, quadrature points, particles, samples, or
design points. Executing OTI arithmetic inside Kokkos kernels enables these
evaluations to run on GPUs while preserving the same derivative semantics used
on the CPU. This supports a path toward GPU-accelerated high-order
sensitivity analysis in larger simulation workflows.

\section{Conclusions}

\texttt{cpp\_oti\_lib} provides a compact header-only C++ implementation of
scalar static OTI arithmetic for automatic differentiation. The library
represents \(\mathbb{OTI}_m^n\) algebras as fixed-size C++ template
instantiations, uses overloaded arithmetic and elementary functions to
propagate coefficients, and exposes derivative values through normalized and
ordinary partial-derivative accessors.

The software includes focused unit tests, generated logs for reporting,
Python bindings for testing and visualization, and illustrative examples that
demonstrate one-, two-, and three-dimensional derivative workflows. The Kokkos
backend preserves the same user-facing arithmetic interface while enabling
OTI derivative propagation on performance-portable CPU and GPU execution
spaces. Together, these components provide a practical route for introducing
higher-order OTI automatic differentiation into modern C++ simulation codes.

\bibliographystyle{unsrtnat}
\bibliography{cpp_oti_lib_documentation}

\end{document}
